\begin{thebibliography}{999}
\bibitem{Lyons2005} R.~Lyons, \emph{Asymptotic enumeration of spanning trees}, Combin.\ Probab.\ Comput.\ \textbf{14} (2005), no.~4, 491--522. [Primary-source verified: Cambridge UP, DOI 10.1017/S096354830500684X.]

\bibitem{Kesten1959} H.~Kesten, \emph{Symmetric random walks on groups}, Trans.\ Amer.\ Math.\ Soc.\ \textbf{92} (1959), no.~2, 336--354. [Crossref-verified: DOI 10.1090/S0002-9947-1959-0109367-6.]

\bibitem{McKay1981} B.~D.~McKay, \emph{The expected eigenvalue distribution of a large regular graph}, Linear Algebra Appl.\ \textbf{40} (1981), 203--216. [Crossref-verified: DOI 10.1016/0024-3795(81)90150-6.]

\bibitem{Williams} R.~F.~Williams, \emph{Classification of subshifts of finite type}, Ann.\ of Math.\ (2) \textbf{98} (1973), no.~1, 120--153; errata ibid.\ \textbf{99} (1974), no.~2, 380--381. [Main paper Crossref-verified (DOI 10.2307/1970908; Crossref records the start page only, the end page 153 is the standard record); errata primary-source verified: JSTOR, DOI 10.2307/1970903.]

\bibitem{Pimsner} M.~V.~Pimsner, \emph{A class of $C^*$-algebras generalizing both Cuntz--Krieger algebras and crossed products by $\Z$}, in: Free probability theory (Waterloo, ON, 1995), Fields Inst.\ Commun.\ \textbf{12}, Amer.\ Math.\ Soc., Providence, RI, 1997, pp.~189--212. [Crossref-verified: DOI 10.1090/fic/012/08, pp.~189--212; Crossref dates the chapter December 1996, the volume is conventionally cited 1997.]

\bibitem{Sunada} T.~Sunada, \emph{Unitary representations of fundamental groups and the spectrum of twisted Laplacians}, Topology \textbf{28} (1989), no.~2, 125--132. [Crossref-verified: DOI 10.1016/0040-9383(89)90015-3.]

\bibitem{Levine} L.~Levine, \emph{Sandpile groups and spanning trees of directed line graphs}, J.\ Combin.\ Theory Ser.\ A \textbf{118} (2011), no.~2, 350--364. [Crossref-verified: DOI 10.1016/j.jcta.2010.04.001.]

\bibitem{Vallieres} D.~Vallières, \emph{On abelian $\ell$-towers of multigraphs}, Ann.\ Math.\ Qu\'ebec \textbf{45} (2021), no.~2, 433--452. [Crossref-verified: DOI 10.1007/s40316-020-00152-4. Both earlier out-of-band page ranges (335--359 and 133--152) were incorrect.]

\bibitem{McGownVallieres} K.~J.~McGown and D.~Vallières, \emph{On abelian $\ell$-towers of multigraphs II}, Ann.\ Math.\ Qu\'ebec \textbf{47}, no.~2, 461--473. [Crossref-verified: DOI 10.1007/s40316-021-00183-5; Crossref carries an online date of November 2021 and an issue date of October 2023 for this record. Both earlier out-of-band venues (J.\ Number Theory \textbf{238} and Ann.\ Math.\ Qu\'ebec \textbf{46}) were incorrect. A third part in the series exists and is not verified here.]

\bibitem{Gonet} S.~R.~Gonet, \emph{Iwasawa theory of Jacobians of graphs}, Algebraic Combinatorics \textbf{5} (2022), no.~5, 827--848. [Crossref-verified: DOI 10.5802/alco.225. This supersedes the earlier entry entirely: the previously listed title (``Jacobians of finite and infinite voltage covers of graphs'') and venue (J.\ Algebraic Combin.) were incorrect in \emph{both} out-of-band passes --- note that Algebraic Combinatorics and J.\ Algebraic Combin.\ are distinct journals.]

\bibitem{CK} J.~Cuntz and W.~Krieger, \emph{A class of $C^*$-algebras and topological Markov chains}, Invent.\ Math.\ \textbf{56} (1980), 251--268; together with the standard graph-algebra $K$-theory and simplicity/pure-infiniteness criteria descending from it.

\bibitem{Powers} R.~T.~Powers, \emph{Representations of uniformly hyperfinite algebras and their associated von Neumann rings}, Ann.\ of Math.\ (2) \textbf{86} (1967), no.~1, 138--171. [Primary-source verified: JSTOR, DOI 10.2307/1970364.]

\bibitem{ArakiWoods} H.~Araki and E.~J.~Woods, \emph{A classification of factors}, Publ.\ Res.\ Inst.\ Math.\ Sci.\ \textbf{4} (1968), no.~1, 51--130. [Primary-source verified: EMS Press / PRIMS article page.]

\bibitem{Krieger} W.~Krieger, \emph{On ergodic flows and the isomorphism of factors}, Math.\ Ann.\ \textbf{223} (1976), 19--70. [Primary-source verified: Springer, Math.\ Ann.\ \textbf{223}, pp.~19--70, February 1976.]

\bibitem{Connes} A.~Connes, \emph{Classification of injective factors. Cases $\mathrm{II}_1$, $\mathrm{II}_\infty$, $\mathrm{III}_\lambda$, $\lambda\neq1$}, Ann.\ of Math.\ (2) \textbf{104} (1976), no.~1, 73--115. [Primary-source verified: JSTOR, DOI 10.2307/1971057.]

\bibitem{FeldmanMoore} J.~Feldman and C.~C.~Moore, \emph{Ergodic equivalence relations, cohomology, and von Neumann algebras.\ I, II}, Trans.\ Amer.\ Math.\ Soc.\ \textbf{234} (1977), 289--324 and 325--359. [Both parts verified: Part I Crossref-verified, DOI 10.1090/S0002-9947-1977-0578656-4, pp.~289--324; Part II primary-source verified, DOI 10.1090/S0002-9947-1977-0578730-2, pp.~325--359. The related 1975 Bulletin announcement, Bull.\ Amer.\ Math.\ Soc.\ \textbf{81} (1975), no.~5, 921--924, DOI 10.1090/S0002-9904-1975-13888-2, is a distinct paper and is not the reference intended here.]

\bibitem{LevineIII} L.~Levine, \emph{Sandpile groups and spanning trees of directed line graphs}, J.\ Combin.\ Theory Ser.\ A \textbf{118} (2011), no.~2, 350--364. [Crossref-verified: DOI 10.1016/j.jcta.2010.04.001.]

\bibitem{Washington} L.~C.~Washington, \emph{Introduction to Cyclotomic Fields}, 2nd ed., Grad.\ Texts in Math.\ \textbf{83}, Springer, 1997; Chapter~13 for the structure theory of $\La$-modules. [Bibliographic data primary-source verified via the Springer book page: author, title, second edition, GTM~\textbf{83}, \copyright~1997. The chapter-level pointer is a reading aid from recollection.]

\bibitem{Ba} H.~Bass, \emph{The Ihara--Selberg zeta function of a tree lattice}, Internat.\ J.\ Math.\ \textbf{3} (1992), 717--797.

\bibitem{GS} R.~Greenberg and G.~Stevens, \emph{$p$-adic $L$-functions and $p$-adic periods of modular forms}, Invent.\ Math.\ \textbf{111} (1993), 407--447. [DOI 10.1007/BF01231294.]

\bibitem{MTT} B.~Mazur, J.~Tate and J.~Teitelbaum, \emph{On $p$-adic analogues of the conjectures of Birch and Swinnerton-Dyer}, Invent.\ Math.\ \textbf{84} (1986), 1--48. [DOI 10.1007/BF01388731.]

\bibitem{ST} H.~M.~Stark and A.~A.~Terras, \emph{Zeta functions of finite graphs and coverings}, Adv.\ Math.\ \textbf{121} (1996), 124--165.

\bibitem{Wa} L.~C.~Washington, \emph{Introduction to Cyclotomic Fields}, 2nd ed., Grad.\ Texts in Math.\ \textbf{83}, Springer, 1997; Chapter~7 for Weierstrass preparation and Chapter~13 for $\La$-modules.

\bibitem{CC} A.~H.~Chamseddine and A.~Connes, \emph{The spectral action principle}, Comm.\ Math.\ Phys.\ \textbf{186} (1997), 731--750.

\bibitem{Co} A.~Connes, \emph{Noncommutative geometry}, Academic Press, San Diego, 1994.

\bibitem{JLO} A.~Jaffe, A.~Lesniewski and K.~Osterwalder, \emph{Quantum $K$-theory. I. The Chern character}, Comm.\ Math.\ Phys.\ \textbf{118} (1988), 1--14.

\bibitem{KS1} M.~Kotani and T.~Sunada, \emph{Albanese maps and off diagonal long time asymptotics for the heat kernel}, Comm.\ Math.\ Phys.\ \textbf{209} (2000), 633--670.

\bibitem{KS2} M.~Kotani and T.~Sunada, \emph{Standard realizations of crystal lattices via harmonic maps}, Trans.\ Amer.\ Math.\ Soc.\ \textbf{353} (2001), 1--20.

\bibitem{Lo} J.-L.~Loday, \emph{Cyclic homology}, 2nd ed., Grundlehren Math.\ Wiss.\ \textbf{301}, Springer, Berlin, 1998.

\bibitem{CFKSV} J.~Coates, T.~Fukaya, K.~Kato, R.~Sujatha and O.~Venjakob, \emph{The $\mathrm{GL}_2$ main conjecture for elliptic curves without complex multiplication}, Publ.\ Math.\ Inst.\ Hautes \'Etudes Sci.\ \textbf{101} (2005), 163--208. [Primary-source verified: the IH\'ES journal page, Volume 101 (2005), pp.~163--208, authors and title as listed.]

\bibitem{Co:XIII} P.~M.~Cohn, \emph{Free rings and their relations}, 2nd ed., London Math.\ Soc.\ Monographs \textbf{19}, Academic Press, London, 1985.

\bibitem{NR} A.~Neeman and A.~Ranicki, \emph{Noncommutative localisation in algebraic $K$-theory I}, Geom.\ Topol.\ \textbf{8} (2004), 1385--1425.

\bibitem{Sch} A.~H.~Schofield, \emph{Representations of rings over skew fields}, London Math.\ Soc.\ Lecture Note Ser.\ \textbf{92}, Cambridge Univ.\ Press, Cambridge, 1985.

\bibitem{KR} E.~Kirchberg and M.~R\o rdam, \emph{Non-simple purely infinite $C^*$-algebras}, Amer.\ J.\ Math.\ \textbf{122} (2000), 637--666.

\bibitem{Ph} N.~C.~Phillips, \emph{A classification theorem for nuclear purely infinite simple $C^*$-algebras}, Doc.\ Math.\ \textbf{5} (2000), 49--114.

\bibitem{R} I.~Raeburn, \emph{Graph algebras}, CBMS Regional Conference Series in Mathematics \textbf{103}, Amer.\ Math.\ Soc., Providence, RI, 2005.

\bibitem{Ro} M.~R\o rdam, \emph{Classification of nuclear, simple $C^*$-algebras}, in: Classification of nuclear $C^*$-algebras. Entropy in operator algebras, Encyclopaedia Math.\ Sci.\ \textbf{126}, Springer, Berlin, 2002, pp.~1--145.

\bibitem{RS} J.~Rosenberg and C.~Schochet, \emph{The K\"unneth theorem and the universal coefficient theorem for Kasparov's generalized $K$-functor}, Duke Math.\ J.\ \textbf{55} (1987), 431--474.

\bibitem{Bla} B.~Blackadar, \emph{$K$-theory for operator algebras}, 2nd ed., Math.\ Sci.\ Res.\ Inst.\ Publ.\ \textbf{5}, Cambridge Univ.\ Press, Cambridge, 1998.

\bibitem{Kas} G.~G.~Kasparov, \emph{The operator $K$-functor and extensions of $C^*$-algebras}, Math.\ USSR-Izv.\ \textbf{16} (1981), no.~3, 513--572.

\bibitem{KP} J.~Kaminker and I.~Putnam, \emph{$K$-theoretic duality for shifts of finite type}, Comm.\ Math.\ Phys.\ \textbf{187} (1997), 509--522.

\bibitem{CMSZ} D.~I.~Cartwright, A.~M.~Mantero, T.~Steger and A.~Zappa, \emph{Groups acting simply transitively on the vertices of a building of type $\widetilde A_2$, I}, Geom.\ Dedicata \textbf{47} (1993), 143--166.

\bibitem{DKM} A.~M.~Duval, C.~J.~Klivans and J.~L.~Martin, \emph{Critical groups of simplicial complexes}, Ann.\ Comb.\ \textbf{17} (2013), 53--70.

\bibitem{KL} M.-H.~Kang and W.-C.~W.~Li, \emph{Zeta functions of complexes arising from $\mathrm{PGL}(3)$}, Adv.\ Math.\ \textbf{256} (2014), 46--103.

\bibitem{KOR} K.~H.~Kim, N.~S.~Ormes and F.~W.~Roush, \emph{The spectra of nonnegative integer matrices via formal power series}, J.\ Amer.\ Math.\ Soc.\ \textbf{13} (2000), 773--806.

\bibitem{BH91} M.~Boyle and D.~Handelman, \emph{The spectra of nonnegative matrices via symbolic dynamics}, Ann.\ of Math.\ (2) \textbf{133} (1991), 249--316.

\bibitem{BH93} M.~Boyle and D.~Handelman, \emph{Algebraic shift equivalence and primitive matrices}, Trans.\ Amer.\ Math.\ Soc.\ \textbf{336} (1993), 121--149.

\bibitem{BS} M.~Boyle and S.~Schmieding, \emph{Strong shift equivalence and the generalized spectral conjecture for nonnegative matrices}, Linear Algebra Appl.\ \textbf{498} (2016), 231--243; arXiv:1501.04697.

\bibitem{BSw} W.~Ballmann and J.~\'Swi\k{a}tkowski, \emph{On $L^{2}$-cohomology and property (T) for automorphism groups of polyhedral cell complexes}, Geom.\ Funct.\ Anal.\ \textbf{7} (1997), 615--645.

\bibitem{Ga} H.~Garland, \emph{$p$-adic curvature and the cohomology of discrete subgroups of $p$-adic groups}, Ann.\ of Math.\ (2) \textbf{97} (1973), no.~3, 375--423. [Primary-source verified: JSTOR record, DOI 10.2307/1970829.]

\bibitem{KLW} M.-H.~Kang, W.-C.~W.~Li and C.-J.~Wang, \emph{The zeta functions of complexes from $\mathrm{PGL}(3)$: a representation-theoretic approach}, Israel J.\ Math.\ \textbf{177} (2010), 335--348.

\bibitem{LSV} A.~Lubotzky, B.~Samuels and U.~Vishne, \emph{Ramanujan complexes of type $\widetilde A_d$}, Israel J.\ Math.\ \textbf{149} (2005), 267--299.

\bibitem{KFH} A.~J.~Koll\'ar, M.~Fitzpatrick, and A.~A.~Houck, \emph{Hyperbolic lattices in circuit quantum electrodynamics}, Nature \textbf{571}, 45 (2019).

\bibitem{MR1} J.~Maciejko and S.~Rayan, \emph{Hyperbolic band theory}, Sci.\ Adv.\ \textbf{7}, eabe9170 (2021).

\bibitem{MR2} J.~Maciejko and S.~Rayan, \emph{Automorphic Bloch theorems for hyperbolic lattices}, Proc.\ Natl.\ Acad.\ Sci.\ USA \textbf{119}, e2116869119 (2022).

\bibitem{Leng22} P.~M.~Lenggenhager \emph{et al.}, \emph{Simulating hyperbolic space on a circuit board}, Nat.\ Commun.\ \textbf{13}, 4373 (2022).

\bibitem{Huang24} L.~Huang \emph{et al.}, \emph{Hyperbolic photonic topological insulators}, Nat.\ Commun.\ \textbf{15}, 1647 (2024).

\bibitem{Xu25} X.~Xu, A.~A.~Mahmoud, N.~Gorgichuk, R.~Thomale, S.~Rayan, and M.~Mariantoni, \emph{A scalable superconducting circuit framework for emulating physics in hyperbolic space}, arXiv:2510.23827 (2025).

\bibitem{KS3} M.~Kotani and T.~Sunada, \emph{Jacobian tori associated with a finite graph and its abelian covering graphs}, Adv.\ Appl.\ Math.\ \textbf{24}, 89 (2000).

\bibitem{BF} M.~Baker and X.~Faber, \emph{Metric properties of the tropical Abel--Jacobi map}, J.\ Algebr.\ Comb.\ \textbf{33}, 349 (2011).

\bibitem{KorS} E.~Korotyaev and N.~Saburova, \emph{Effective masses for Laplacians on periodic graphs}, J.\ Math.\ Anal.\ Appl.\ \textbf{436}, 104 (2016).

\bibitem{LSV:Phys1} A.~Lubotzky, B.~Samuels, and U.~Vishne, \emph{Ramanujan complexes of type $\widetilde A_d$}, Isr.\ J.\ Math.\ \textbf{149}, 267 (2005).

\bibitem{ANCFT9} R.~Chen, \emph{Algorithmic non-commutative class field theory, IX: exceptional zeros, the graph $\mathcal L$-invariant, and the operator-algebraic Mazur--Tate--Teitelbaum formula}, preprint (2026).

\bibitem{ANCFT10} R.~Chen, \emph{Algorithmic non-commutative class field theory, X: $\widetilde A_2$ complexes, higher simplicial Laplacians, and the fate of the trace obstruction in rank two}, preprint (2026).

\bibitem{ANCFT12} R.~Chen, \emph{Algorithmic non-commutative class field theory, XII: $\widetilde A_2$ towers, no abelian sector, the covering norm theorem, and the two digraph towers}, preprint (2026).

\bibitem{ANCFT14} R.~Chen, \emph{Algorithmic non-commutative class field theory, XIV: the graph $\mathcal L$-invariant as spectral action, Bloch effective mass and Albanese length, and what the cyclic Chern character sees}, preprint (2026).

\bibitem{BHN} R.~Bacher, P.~de~la~Harpe and T.~Nagnibeda, \emph{The lattice of integral flows and the lattice of integral cuts on a finite graph}, Bull.\ Soc.\ Math.\ France \textbf{125} (1997), 167--198.

\bibitem{Biggs} N.~Biggs, \emph{Chip-firing and the critical group of a graph}, J.\ Algebraic Combin.\ \textbf{9} (1999), 25--45.

\bibitem{Bowen14} L.~Bowen, \emph{The type and stable type of the boundary of a Gromov hyperbolic group}, Geom.\ Dedicata \textbf{172} (2014), 363--386.

\bibitem{BDF87} D.~Buchholz, C.~D'Antoni and K.~Fredenhagen, \emph{The universal structure of local algebras}, Comm.\ Math.\ Phys.\ \textbf{111} (1987), 123--135.

\bibitem{CPW23} V.~Chandrasekaran, G.~Penington and E.~Witten, \emph{Large $N$ algebras and generalized entropy}, JHEP \textbf{04} (2023), 009.

\bibitem{CLPW23} V.~Chandrasekaran, R.~Longo, G.~Penington and E.~Witten, \emph{An algebra of observables for de Sitter space}, JHEP \textbf{02} (2023), 082.

\bibitem{ANCFT} R.~Chen, \emph{Algorithmic non-commutative class field theory, I--XVI}, preprints (2026); here III (the line-digraph recursion and the Euler term), IV (the limit module), IX (the defects), X--XII (rank two).

\bibitem{CLM08} G.~Cornelissen, O.~Lorscheid and M.~Marcolli, \emph{On the $K$-theory of graph $C^*$-algebras}, Acta Appl.\ Math.\ \textbf{102} (2008), 57--69.

\bibitem{CM13} G.~Cornelissen and M.~Marcolli, \emph{Graph reconstruction and quantum statistical mechanics}, J.\ Geom.\ Phys.\ \textbf{72} (2013), 110--117.

\bibitem{Cuntz81} J.~Cuntz, \emph{A class of $C^*$-algebras and topological Markov chains II: reducible chains and the Ext-functor for $C^*$-algebras}, Invent.\ Math.\ \textbf{63} (1981), 23--50.

\bibitem{CR03} P.~Cutting and G.~Robertson, \emph{Type III actions on boundaries of $\widetilde A_n$ buildings}, J.\ Operator Theory \textbf{49} (2003), 25--44.

\bibitem{FLMO23} K.~Furuya, N.~Lashkari, M.~Moosa and S.~Ouseph, \emph{Information loss, mixing and emergent type $\mathrm{III}_1$ factors}, JHEP \textbf{08} (2023), 111.

\bibitem{GMP22} E.~Gesteau, M.~Marcolli and S.~Parikh, \emph{Holographic tensor networks from hyperbolic buildings}, JHEP \textbf{10} (2022), 169.

\bibitem{GLS22} P.~Gorantla, H.~T.~Lam and S.-H.~Shao, \emph{Fractons on graphs and complexity}, Phys.\ Rev.\ B \textbf{106} (2022), 195139.

\bibitem{GLSS23} P.~Gorantla, H.~T.~Lam, N.~Seiberg and S.-H.~Shao, \emph{Gapped lineon and fracton models on graphs}, Phys.\ Rev.\ B \textbf{107} (2023), 125121.

\bibitem{GKPSW} S.~S.~Gubser, J.~Knaute, S.~Parikh, A.~Samberg and P.~Witaszczyk, \emph{$p$-adic AdS/CFT}, Comm.\ Math.\ Phys.\ \textbf{352} (2017), 1019--1059.

\bibitem{GubserParikh} S.~S.~Gubser and S.~Parikh, \emph{Geodesic bulk diagrams on the Bruhat--Tits tree}, Phys.\ Rev.\ D \textbf{96} (2017), 066024.

\bibitem{HMSS} M.~Heydeman, M.~Marcolli, I.~Saberi and B.~Stoica, \emph{Tensor networks, $p$-adic fields, and algebraic curves: arithmetic and the $\mathrm{AdS}_3/\mathrm{CFT}_2$ correspondence}, Adv.\ Theor.\ Math.\ Phys.\ \textbf{22} (2018), 93--176.

\bibitem{HuangJepsen} A.~Huang and C.~B.~Jepsen, \emph{Finite temperature at finite places}, JHEP \textbf{03} (2025), 097.

\bibitem{INO08} M.~Izumi, S.~Neshveyev and R.~Okayasu, \emph{The ratio set of the harmonic measure of a random walk on a hyperbolic group}, Israel J.\ Math.\ \textbf{163} (2008), 285--316.

\bibitem{Krieger76} W.~Krieger, \emph{On ergodic flows and the isomorphism of factors}, Math.\ Ann.\ \textbf{223} (1976), 19--70.

\bibitem{LL1} S.~Leutheusser and H.~Liu, \emph{Causal connectability between quantum systems and the black hole interior in holographic duality}, Phys.\ Rev.\ D \textbf{108} (2023), 086019.

\bibitem{LL2} S.~Leutheusser and H.~Liu, \emph{Emergent times in holographic duality}, Phys.\ Rev.\ D \textbf{108} (2023), 086020.

\bibitem{Marcolli20} M.~Marcolli, \emph{Holographic codes on Bruhat--Tits buildings and Drinfeld symmetric spaces}, Pure Appl.\ Math.\ Q.\ \textbf{16} (2020), 1--33.

\bibitem{Mondal} A.~Mondal, S.~Parikh, P.~Pradhan and R.~Sengar, \emph{Toward holography on biregular trees}, Phys.\ Rev.\ D \textbf{112} (2025), 086011.

\bibitem{RR97} J.~Ramagge and G.~Robertson, \emph{Factors from trees}, Proc.\ Amer.\ Math.\ Soc.\ \textbf{125} (1997), 2051--2055.

\bibitem{Rob05} G.~Robertson, \emph{Boundary operator algebras for free uniform tree lattices}, Houston J.\ Math.\ \textbf{31} (2005), 913--935.

\bibitem{Wiesbrock93} H.-W.~Wiesbrock, \emph{Half-sided modular inclusions of von Neumann algebras}, Comm.\ Math.\ Phys.\ \textbf{157} (1993), 83--92.

\bibitem{Witten22} E.~Witten, \emph{Gravity and the crossed product}, JHEP \textbf{10} (2022), 008.

\bibitem{BH09} J.~S.~Baras and P.~Hovareshti, \emph{Efficient and robust communication topologies for distributed decision making in networked systems}, Proc.\ 48th IEEE Conf.\ Decision and Control (2009).

\bibitem{BK11} H.~Bidkhori and S.~Kishore, \emph{Counting the spanning trees of a directed line graph}, Combin.\ Probab.\ Comput.\ \textbf{20} (2011), 11--25.

\bibitem{Boesch86} F.~T.~Boesch, \emph{On unreliability polynomials and graph connectivity in reliable network synthesis}, J.\ Graph Theory \textbf{10} (1986), 339--352.

\bibitem{CHP15} S.~H.~Chan, H.~D.~L.~Hollmann and D.~V.~Pasechnik, \emph{Sandpile groups of generalized de Bruijn and Kautz graphs and circulant matrices over finite fields}, J.\ Algebra \textbf{421} (2015), 268--295.

\bibitem{Colbourn87} C.~J.~Colbourn, \emph{The Combinatorics of Network Reliability}, Oxford University Press, 1987.

\bibitem{CM81} A.~A.~Cuoco and P.~Monsky, \emph{Class numbers in $\Z_p^{d}$-extensions}, Math.\ Ann.\ \textbf{255} (1981), 235--258.

\bibitem{DV22} S.~DuBose and D.~Vall\`ieres, \emph{On $\Z_\ell^{d}$-towers of graphs}, preprint, arXiv:2207.01711 (2022).

\bibitem{FAY83} M.~A.~Fiol, I.~Alegre and J.~L.~A.~Yebra, \emph{Line digraph iterations and the $(d,k)$ problem for directed graphs}, Proc.\ 10th Int.\ Symp.\ Computer Architecture (1983), 174--177.

\bibitem{FYA84} M.~A.~Fiol, J.~L.~A.~Yebra and I.~Alegre, \emph{Line digraph iterations and the $(d,k)$ digraph problem}, IEEE Trans.\ Comput.\ \textbf{C-33} (1984), 400--403.

\bibitem{FL92} M.~A.~Fiol and A.~S.~Llad\'o, \emph{The partial line digraph technique in the design of large interconnection networks}, IEEE Trans.\ Comput.\ \textbf{41} (1992), 848--857.

\bibitem{Gonet22} S.~R.~Gonet, \emph{Iwasawa theory of Jacobians of graphs}, Algebraic Combinatorics \textbf{5} (2022), 827--848.

\bibitem{Kataoka26} T.~Kataoka, \emph{Construction of graph coverings with prescribed Iwasawa invariants}, preprint, arXiv:2603.23088 (2026).

\bibitem{KM24} S.~Kleine and K.~M\"uller, \emph{On the growth of the Jacobian in $\Z_p^{l}$-voltage covers of graphs}, Algebraic Combinatorics \textbf{7} (2024), no.~4.

\bibitem{Knuth67} D.~E.~Knuth, \emph{Oriented subtrees of an arc digraph}, J.\ Combin.\ Theory \textbf{3} (1967), 309--314.

\bibitem{Levine11} L.~Levine, \emph{Sandpile groups and spanning trees of directed line graphs}, J.\ Combin.\ Theory Ser.~A \textbf{118} (2011), 350--364.

\bibitem{MV24} K.~McGown and D.~Vall\`ieres, \emph{On abelian $\ell$-towers of multigraphs III}, Ann.\ Math.\ Qu\'ebec \textbf{48} (2024), 1--19.

\bibitem{SHM05} C.~Song, S.~Havlin and H.~A.~Makse, \emph{Self-similarity of complex networks}, Nature \textbf{433} (2005), 392--395.

\bibitem{SHM06} C.~Song, S.~Havlin and H.~A.~Makse, \emph{Origins of fractality in the growth of complex networks}, Nature Physics \textbf{2} (2006), 275--281.

\bibitem{Val21} D.~Vall\`ieres, \emph{On abelian $\ell$-towers of multigraphs}, Ann.\ Math.\ Qu\'ebec \textbf{45} (2021), 433--452.

\bibitem{ADFB96} B.~D.~O.~Anderson, M.~Deistler, L.~Farina and L.~Benvenuti, \emph{Nonnegative realization of a linear system with nonnegative impulse response}, IEEE Trans.\ Circuits Syst.\ I \textbf{43} (1996), 134--142.

\bibitem{Ben20ELA} L.~Benvenuti, \emph{The NIEP and the positive realization problem}, Electron.\ J.\ Linear Algebra \textbf{36} (2020), 367--384.

\bibitem{Ben20SCL} L.~Benvenuti, \emph{A lower bound on the dimension of minimal positive realizations for discrete time systems}, Systems Control Lett.\ \textbf{135} (2020), 104595.

\bibitem{BF99} L.~Benvenuti and L.~Farina, \emph{An example of how positivity may force realizations of ``large'' dimension}, Systems Control Lett.\ \textbf{36} (1999), 261--266.

\bibitem{BF04} L.~Benvenuti and L.~Farina, \emph{A tutorial on the positive realization problem}, IEEE Trans.\ Automat.\ Control \textbf{49} (2004), 651--664.

\bibitem{Farina96} L.~Farina, \emph{On the existence of a positive realization}, Systems Control Lett.\ \textbf{28} (1996), 219--226.

\bibitem{KOR00} K.~H.~Kim, N.~S.~Ormes and F.~W.~Roush, \emph{The spectra of nonnegative integer matrices via formal power series}, J.\ Amer.\ Math.\ Soc.\ \textbf{13} (2000), 773--806.

\bibitem{LM95} D.~Lind and B.~Marcus, \emph{An Introduction to Symbolic Dynamics and Coding}, Cambridge University Press, 1995.

\bibitem{LoewyLondon78} R.~Loewy and D.~London, \emph{A note on an inverse problem for nonnegative matrices}, Linear Multilinear Algebra \textbf{6} (1978), 83--90.

\bibitem{NM03} B.~Nagy and M.~Matolcsi, \emph{A lowerbound on the dimension of positive realizations}, IEEE Trans.\ Circuits Syst.\ I \textbf{50} (2003), 782--784.

\bibitem{OMK84} Y.~Ohta, H.~Maeda and S.~Kodama, \emph{Reachability, observability and realizability of continuous-time positive systems}, SIAM J.\ Control Optim.\ \textbf{22} (1984), 171--180.

\bibitem{TS25} H.~Taghavian and J.~Sj\"olund, \emph{Minimal positive Markov realizations}, preprint, arXiv:2502.21102 (2025).

\bibitem{AEdB} T.~van Aardenne-Ehrenfest and N.~G.~de Bruijn, \emph{Circuits and trees in oriented linear graphs}, Simon Stevin \textbf{28} (1951), 203--217.

\bibitem{BBT} G.~Bresler, M.~Bresler and D.~Tse, \emph{Optimal assembly for high throughput shotgun sequencing}, BMC Bioinformatics \textbf{14} (Suppl.~5) (2013), S18.

\bibitem{BidkhoriKishore} H.~Bidkhori and S.~Kishore, \emph{Counting the spanning trees of a directed line graph}, Combin.\ Probab.\ Comput.\ \textbf{20} (2011), 11--25.

\bibitem{CHP} S.~H.~Chan, H.~D.~L.~Hollmann and D.~V.~Pasechnik, \emph{Sandpile groups of generalized de Bruijn and Kautz graphs and circulant matrices over finite fields}, J.\ Algebra \textbf{421} (2015), 268--295.

\bibitem{Bio2} R.~Chen, \emph{The sandpile group as a coordinate system on genome reconstructions}, companion preprint (Bio~2), 2026.

\bibitem{KSP} C.~Kingsford, M.~C.~Schatz and M.~Pop, \emph{Assembly complexity of prokaryotic genomes using short reads}, BMC Bioinformatics \textbf{11} (2010), 21.

\bibitem{Knuth} D.~E.~Knuth, \emph{Oriented subtrees of an arc digraph}, J.\ Combin.\ Theory \textbf{3} (1967), 309--314.

\bibitem{Loukides} G.~Loukides, S.~P.~Pissis et al., \emph{Beyond the BEST theorem: fast assessment of Eulerian trails}, in: Fundamentals of Computation Theory (FCT 2021), Lecture Notes in Comput.\ Sci.\ \textbf{12867}, Springer, 2021, pp.~162--175.

\bibitem{NagarajanPop} N.~Nagarajan and M.~Pop, \emph{Parametric complexity of sequence assembly: theory and applications to next generation sequencing}, J.\ Comput.\ Biol.\ \textbf{16} (2009), 897--908.

\bibitem{ObscuraTomescu} N.~Obscura Acosta and A.~I.~Tomescu, \emph{Simplicity in Eulerian circuits: uniqueness and safety}, preprint, arXiv:2208.08522 (2022).

\bibitem{RT} V.~Reiner and D.~Tseng, \emph{Critical groups of covering, voltage and signed graphs}, Discrete Math.\ \textbf{318} (2014), 10--40.

\bibitem{SmithTutte} C.~A.~B.~Smith and W.~T.~Tutte, \emph{On unicursal paths in a network of degree 4}, Amer.\ Math.\ Monthly \textbf{48} (1941), 233--237.

\bibitem{Zaslavsky} T.~Zaslavsky, \emph{Signed graphs}, Discrete Appl.\ Math.\ \textbf{4} (1982), 47--74.

\bibitem{AE} S.~F.~Altschul and B.~W.~Erickson, \emph{Significance of nucleotide sequence alignments: a method for random sequence permutation that preserves dinucleotide and codon usage}, Mol.\ Biol.\ Evol.\ \textbf{2} (1985), 526--538.

\bibitem{ABCS} R.~Arratia, B.~Bollob\'as, D.~Coppersmith and G.~B.~Sorkin, \emph{Euler circuits and DNA sequencing by hybridization}, Discrete Appl.\ Math.\ \textbf{104} (2000), 63--96.

\bibitem{BK} H.~Bidkhori and S.~Kishore, \emph{Counting the spanning trees of a directed line graph}, Combin.\ Probab.\ Comput.\ \textbf{20} (2011), 11--25.

\bibitem{BW} G.~R.~Brightwell and P.~Winkler, \emph{Counting Eulerian circuits is $\#$P-complete}, in: Proc.\ ALENEX/ANALCO 2005, SIAM, 2005, pp.~259--262.

\bibitem{Bio1} R.~Chen, \emph{Assembly ambiguity is not a function of branch statistics: the sandpile group of a de Bruijn graph}, companion preprint (Bio~1, version~2), 2026.

\bibitem{CMN} C.~J.~Colbourn, W.~J.~Myrvold and E.~Neufeld, \emph{Two algorithms for unranking arborescences}, J.\ Algorithms \textbf{20} (1996), 268--281.

\bibitem{CC:Bio2} P.~Creed and M.~Cryan, \emph{The number of Euler tours of random directed graphs}, Electron.\ J.\ Combin.\ \textbf{20} (2013), P13.

\bibitem{HLMPPW} A.~E.~Holroyd, L.~Levine, K.~M\'esz\'aros, Y.~Peres, J.~Propp and D.~B.~Wilson, \emph{Chip-firing and rotor-routing on directed graphs}, in: In and Out of Equilibrium 2, Progress in Probability \textbf{60}, Birkh\"auser, 2008, pp.~331--364.

\bibitem{uShuffle} M.~Jiang, J.~Anderson, J.~Gillespie and M.~Mayne, \emph{uShuffle: a useful tool for shuffling biological sequences while preserving the $k$-let counts}, BMC Bioinformatics \textbf{9} (2008), 192.

\bibitem{KMUW} D.~Kandel, Y.~Matias, R.~Unger and P.~Winkler, \emph{Shuffling biological sequences}, Discrete Appl.\ Math.\ \textbf{71} (1996), 171--185.

\bibitem{KM} J.~D.~Kececioglu and E.~W.~Myers, \emph{Combinatorial algorithms for DNA sequence assembly}, Algorithmica \textbf{13} (1995), 7--51.

\bibitem{MB} P.~Medvedev and M.~Brudno, \emph{Maximum likelihood genome assembly}, J.\ Comput.\ Biol.\ \textbf{16} (2009), 1101--1116.

\bibitem{PDDK} V.~B.~Priezzhev, D.~Dhar, A.~Dhar and S.~Krishnamurthy, \emph{Eulerian walkers as a model of self-organized criticality}, Phys.\ Rev.\ Lett.\ \textbf{77} (1996), 5079--5082.

\bibitem{PW} J.~G.~Propp and D.~B.~Wilson, \emph{How to get a perfectly random sample from a generic Markov chain and generate a random spanning tree of a directed graph}, J.\ Algorithms \textbf{27} (1998), 170--217.

\bibitem{Wilson} D.~B.~Wilson, \emph{Generating random spanning trees more quickly than the cover time}, in: Proc.\ 28th ACM Symp.\ Theory of Computing (1996), pp.~296--303.

\bibitem{Alar} D.~Alar, M.~Celaya, L.~D.~Garc\'ia Puente, S.~Henson and A.~Wheeler, \emph{The sandpile group of a thick cycle graph}, Electron.\ J.\ Graph Theory Appl.\ \textbf{10} (2022), no.~2.

\bibitem{ChenSchedler} W.~Chen and T.~Schedler, \emph{Concrete and abstract structure of the sandpile group for thick trees with loops}, preprint, arXiv:math/0701381 (2007).

\bibitem{BCALM} R.~Chikhi, A.~Limasset and P.~Medvedev, \emph{Compacting de Bruijn graphs from sequencing data quickly and in low memory}, Bioinformatics \textbf{32} (2016), i201--i208.

\bibitem{Flye} M.~Kolmogorov, J.~Yuan, Y.~Lin and P.~A.~Pevzner, \emph{Assembly of long, error-prone reads using repeat graphs}, Nat.\ Biotechnol.\ \textbf{37} (2019), 540--546.

\bibitem{ABS} R.~Arratia, B.~Bollob\'as and G.~B.~Sorkin, \emph{The interlace polynomial of a graph}, J.\ Combin.\ Theory Ser.~B \textbf{92} (2004), 199--233.

\bibitem{LasVergnas} M.~Las Vergnas, \emph{Le polyn\^ome de Martin d'un graphe eul\'erien}, Ann.\ Discrete Math.\ \textbf{17} (1983), 397--411.

\bibitem{MMN} C.~Merino, I.~Moffatt and S.~Noble, \emph{The critical group of a combinatorial map}, Combinatorial Theory \textbf{5} (2025), no.~3; arXiv:2308.13342.

\bibitem{MacrisPule} N.~Macris and J.~V.~Pul\'e, \emph{An alternative formula for the number of Euler trails for a class of digraphs}, Discrete Math.\ \textbf{154} (1996), 301--305.

\bibitem{Jonsson} J.~Jonsson, \emph{On the number of Euler trails in directed graphs}, Math.\ Scand.\ \textbf{90} (2002), 191--214.

\bibitem{PSX} D.~Perkinson, N.~Salter and T.~Xu, \emph{A note on the critical group of a line graph}, Electron.\ J.\ Combin.\ \textbf{18} (2011), P124.

\bibitem{PTT} P.~A.~Pevzner, H.~Tang and G.~Tesler, \emph{De novo repeat classification and fragment assembly}, Genome Res.\ \textbf{14} (2004), 1786--1796.

\bibitem{PTW} P.~A.~Pevzner, H.~Tang and M.~S.~Waterman, \emph{An Eulerian path approach to DNA fragment assembly}, Proc.\ Natl.\ Acad.\ Sci.\ USA \textbf{98} (2001), 9748--9753.

\bibitem{Bio3} R.~Chen, \emph{The sandpile group of a de Bruijn graph splits along its repeats}, companion preprint (Bio~3), 2026.

\bibitem{Lauri} J.~Lauri, \emph{On a determinant formula for enumerating Euler trails in a class of digraphs}, Rend.\ Sem.\ Mat.\ Messina Ser.~II \textbf{5} (1999), 75--90.

\bibitem{Toth} L.~T\'othm\'er\'esz, \emph{On canonical sandpile actions of embedded graphs}, Advances in Combinatorics (2026), no.~6; arXiv:2504.05275.

\end{thebibliography}
